\section{Ba bài đứng quanh bài này}
\label{sec:ch05-cau-noi}

\index{encoder-decoder!ba bài liên quan}%
Bài 2014 không mọc ra từ đất trống, và ba bài dưới đây là ba lý do khác nhau để
biết chuyện đó. Bài thứ nhất là bài mà chính Sutskever nhường quyền ưu tiên.
Bài thứ hai dựng cùng kiến trúc gần như cùng lúc, ở một chỗ khác, với một chi
tiết khác hẳn. Bài thứ ba là bài đo được thứ mà người ta vẫn gán nhầm cho
chương này. Phân biệt nó với bài thứ hai là chỗ cả chương này lẫn
chương~\ref{ch:dong-chinh} phải cẩn thận nhất.

\begin{bridge}{Kalchbrenner và Blunsom 2013: ai làm trước}
Mục 4 của bài 2014 viết thẳng: công trình của họ gần nhất với Kalchbrenner và
Blunsom, hai người đầu tiên ánh xạ cả câu input thành một vector rồi bung ngược
vector ấy trở lại thành câu (\enquote{who were the first to map the input
sentence into a vector and then back to a sentence}). Chỗ khác nhau, cũng theo
lời họ, là bài 2013 mã hoá câu bằng mạng tích chập, mà mạng tích chập thì làm
mất thứ tự từ (\enquote{lose the ordering of the words}).

Nghĩa là ý tưởng \enquote{nén cả câu vào một vector rồi bung ra} có trước bài
này một năm. Cái bài 2014 thêm vào là dùng một mạng hồi quy có cổng cho cả hai
đầu, nhờ đó thứ tự từ sống sót qua chỗ nén.

EMNLP 2013, trang 1700-1709~\autocite{kalchbrenner2013recurrent}. Tôi chỉ tra
thư mục bài này, chưa đọc nó, nên đây là đúng câu gán quyền ưu tiên của
Sutskever và không hơn.
\end{bridge}

\begin{bridge}{Cho và cộng sự 2014a: cùng kiến trúc, khác một chỗ nối}
Đây là bài chương~\ref{ch:lstm} đã trích cho đơn vị có cổng hai cổng. Nửa còn
lại của nó, và là nửa mang tên bài, là RNN Encoder-Decoder: cùng ý tưởng nén
câu vào một vector, công bố cùng năm.

Chỗ khác nhau nằm ở một chi tiết mà mục~\ref{sec:ch05-mot-vector} đã nhấn. Ở
bài 2014 của Sutskever, \tn{vector ngữ cảnh}{context vector} vào đúng một lần,
làm trạng thái khởi tạo của decoder, rồi thôi. Ở Cho, decoder đọc lại vector ấy
ở \emph{mọi} bước. Hai cách đều là \enquote{encoder-decoder} và chúng đặt chỗ
thắt ở hai chỗ khác nhau: cách thứ nhất bắt vector đi qua toàn bộ chuỗi trạng
thái của decoder, cách thứ hai không.

Mục tiêu của hai bài cũng khác. Cho và cộng sự chủ yếu nhắm vào việc lắp mạng
của họ vào một hệ dịch thống kê có sẵn, tức chấm điểm lại các cụm từ, chứ không
dịch thẳng; chính Sutskever ghi nhận chỗ đó trong mục 4.

EMNLP 2014, trang 1724-1734~\autocite{cho2014encdec}.
\end{bridge}

\begin{bridge}{Cho và cộng sự 2014b: bài thật sự đo độ tụt theo độ dài câu}
Đây là bài mà cả chương này lẫn chương~\ref{ch:dong-chinh} phải gọi đúng tên,
vì gọi nhầm là cách nhanh nhất để kể sai lịch sử.

\enquote{Cho và cộng sự 2014} là \emph{hai} bài khác nhau, và Bahdanau phân biệt
chúng bằng chữ cái: 2014a là bài EMNLP ở hộp trên, còn 2014b là
\enquote{On the Properties of Neural Machine Translation}, một bài workshop
ngắn hơn ở SSST-8, Doha, tháng 10 năm 2014.

Bài 2014b mới là bài đo. Tóm tắt của nó: \enquote{We show that the neural
machine translation performs relatively well on short sentences without unknown
words, but its performance degrades rapidly as the length of the sentence and
the number of unknown words increase}. Và phần mở đầu của Bahdanau dẫn đúng bài
này khi nói \enquote{Cho et al. (2014b) showed that indeed the performance of a
basic encoder-decoder deteriorates rapidly as the length of an input sentence
increases}.

Chú ý luôn vế thứ hai trong câu của bài 2014b: độ dài \emph{và} số từ ngoài từ
điển. Corpus của mục~\ref{sec:ch05-corpus} có từ điển đóng nên không nói được gì
về vế sau.

SSST-8, trang 103-111~\autocite{cho2014properties}. Tôi mới đọc tóm tắt bài này,
chưa đọc thân bài, nên chương chỉ dùng nó ở mức phát biểu và không trích số nào
của nó.
\end{bridge}

\subsection*{Vậy ai nói encoder-decoder hỏng ở câu dài}

\index{encoder-decoder!ai đo độ tụt theo độ dài}%
Cộng ba hộp trên với chính lời bài 2014 thì trả lời được câu hỏi mà cả chương
này xoay quanh, và câu trả lời không phải câu hay được kể.

Sutskever và cộng sự \emph{không} báo cáo mô hình của họ tụt theo độ dài câu.
Ngược lại, và ở năm chỗ.

Ngay trong tóm tắt: \enquote{Additionally, the LSTM did not have difficulty on
long sentences.} Phần mở đầu, mục 1: \enquote{Surprisingly, the LSTM did not
suffer on very long sentences, despite the recent experience of other
researchers with related architectures}. Mục 3.7 của họ tên là
\enquote{Performance on long sentences}
và mở đầu bằng \enquote{We were surprised to discover that the LSTM did well on
long sentences}. Chú thích hình 3 ghi \enquote{There is no degradation on
sentences with less than 35 words, there is only a minor degradation on the
longest sentences}. Có suy giảm, ở nhóm câu dài nhất, và bài gọi nó là nhỏ; đây
là nửa câu hay bị bỏ khi người ta trích chú thích ấy. Kết luận thì nhắc lại lần
nữa, kèm chỗ họ nói rõ đã tưởng ngược lại: \enquote{We were
initially convinced that the LSTM would fail on long sentences due to its
limited memory, and other researchers reported poor performance on long
sentences with a model similar to ours}.

Người đo được độ tụt là Cho và cộng sự trong bài 2014b, trên một mô hình khác,
với chỗ nối khác và bề rộng khác.

Chỗ này còn một lớp nữa, và tôi suýt đọc lướt qua nó. Sutskever cũng nhắc tới
chuyện đó trong mục 4, khi mô tả bài Bahdanau như một bài dùng attention để
khắc phục chỗ mô hình của Cho làm kém trên câu dài (\enquote{to overcome the
poor performance on long sentences experienced by Cho et al.}). Nhưng số trong
ngoặc vuông ngay sau chữ \enquote{Cho et al.} ở bài Sutskever trỏ tới
arXiv:1406.1078, tức bài \emph{2014a}. Sutskever không trích bài 2014b ở đâu
cả.

Hai chuyện ấy không đá nhau, và phân biệt được chúng là phân biệt giữa một mô
hình với một phép đo trên mô hình đó: 2014a là bài dựng ra kiến trúc, 2014b là
bài đem chính kiến trúc ấy ra đo theo độ dài câu. Nhưng nó có nghĩa là câu gán
kết quả đo cho 2014b là câu của Bahdanau, không phải của Sutskever, và chương
này không được mượn Sutskever để chống lưng cho nó.

Năm chỗ, cùng một bài, cùng một hướng. Cái sai không nằm ở chỗ tin rằng
\tn{chỗ thắt}{bottleneck}
có thật, vì mục~\ref{sec:ch05-bop-vector} vừa đo nó. Nó nằm ở chỗ gán kết quả
cho bài không đưa ra kết quả ấy, và một sinh viên mở bài 1409.3215 ra tìm hình
BLEU tụt theo độ dài sẽ tìm mãi không thấy.
